\documentclass[11pt]{article}
\usepackage[margin=1in]{geometry}
\usepackage{amsmath,amssymb,amsthm,mathtools}
\usepackage{graphicx}
\usepackage{microtype}
\usepackage[hidelinks]{hyperref}
\usepackage{enumitem}

\newcommand{\R}{\mathbb R}
\newcommand{\p}{\partial}
\newcommand{\WTP}{\mathrm{Weak\ C6}}
\newtheorem{theorem}{Theorem}

\title{A mixed-derivative obstruction to Weak C6\\
for the three-dimensional gradient}
\author{Candidate counterexample for arXiv:2412.09987, Question 6.4}
\date{19 July 2026}

\begin{document}
\maketitle

\begin{center}
\fbox{\parbox{0.91\textwidth}{
\textbf{Status: candidate counterexample, likely valid.}
The displayed representation in Question 6.4 does not admit property
\(\WTP\).  This is a full negative answer to that precise question, subject
to human review; it is not a solution of the section's broader representation
problem.}}
\end{center}

\section{Source and original question}

The source is Wen Qi Zhang, \emph{Stein-Weiss, and power weight Korn type
Hardy-Sobolev Inequalities in \(L^1\) norm}, arXiv:2412.09987v2, Section 6,
Question 6.4, page 18 \cite{Zhang}.  The question asks whether the following
specific representation of \(\nabla\) on \(\R^3\) admits property
\(\WTP\):
\begin{equation}\label{eq:LK}
L(D)=\begin{bmatrix}
\p_{y_2}\p_{y_3}&\p_{y_1}\p_{y_3}&-2\p_{y_1}\p_{y_2}
\end{bmatrix},\qquad
K(y)=\begin{pmatrix}y_2y_3\\y_1y_3\\-\frac12y_1y_2\end{pmatrix},
\end{equation}
with Green function
\begin{equation}\label{eq:G}
G(x)=\nabla_x\Phi(x),\qquad \Phi(x)=\frac{1}{4\pi|x|}.
\end{equation}
The source then proposes three cubic columns and asks whether they can be used
to construct the required factorization.

\begin{figure}[ht]
\centering
\includegraphics[width=\textwidth]{figures/open_problem_crop.png}
\caption{Question 6.4 and the proposed construction, from page 18 of the
source paper.  The crop retains the full page width and the complete question.}
\label{fig:question}
\end{figure}

\section{Weak C6 in this representation}

Here \(V=F=\R\) and \(E=\R^3\).  The source definition asks for a
finite-dimensional space \(M\) and homogeneous maps
\(T(x):M\to\R\), \(P(y):\R\to M\), independent of the multiindex, such
that
\begin{equation}\label{eq:weak-c6}
\sum_{|\alpha|=2}\p_y^\alpha\bigl(T(x)P(y)\bigr)L_\alpha
=\sum_{|\alpha|=2}\sum_{e_i\leq\alpha}
\binom{\alpha}{e_i}\p_{x_i}G(x)\,
\p_y^{\alpha-e_i}K(y)L_\alpha
\end{equation}
as maps \(E\to V\), for every \(x\ne0\).

The only nonzero coefficient rows of \(L\) are
\begin{equation}\label{eq:Lrows}
L_{23}=(1,0,0),\qquad L_{13}=(0,1,0),\qquad
L_{12}=(0,0,-2).
\end{equation}
For fixed \(x\ne0\), write
\[
Q_x(y):=T(x)P(y),\qquad
H(x):=D^2\Phi(x),\qquad g_i(x):=\p_{x_i}G(x).
\]
Thus \(g_i=(H_{i1},H_{i2},H_{i3})\).

\section{Proof intuition}

Equation \eqref{eq:weak-c6} prescribes three mixed second derivatives of the
\emph{same scalar function} \(Q_x\).  Therefore the first two prescriptions
must agree after one more derivative: both become \(\p_{123}Q_x\).  The
right-hand sides imposed by \eqref{eq:LK}, however, demand incompatible values.
At \(x=(1,0,0)\), the coefficient seen by \(\p_1\p_{23}Q_x\) is
\(-1/(8\pi)\), while the coefficient seen by \(\p_2\p_{13}Q_x\) is
\(5/(8\pi)\).  No choice of the auxiliary space or factorization can alter
this compatibility requirement.

\section{The negative answer}

\begin{theorem}\label{thm:main}
The representation \eqref{eq:LK}--\eqref{eq:G} of \(\nabla\) on \(\R^3\)
does not admit property \(\WTP\).  Equivalently, there is no
finite-dimensional \(M\) and no pair of maps \(T,P\) satisfying the source's
Weak C6 identity for every \(x\in\R^3\setminus\{0\}\).
\end{theorem}

\begin{proof}
Assume such \(M,T,P\) exist, and fix \(x\ne0\).  Since the three rows in
\eqref{eq:Lrows} occupy different coordinates, equality of the row maps in
\eqref{eq:weak-c6} is equivalent, component by component, to
\begin{align}
\p_{23}Q_x&=R_{23}:=g_2\p_3K+g_3\p_2K,\label{eq:R23def}\\
\p_{13}Q_x&=R_{13}:=g_1\p_3K+g_3\p_1K,\label{eq:R13def}\\
\p_{12}Q_x&=R_{12}:=g_1\p_2K+g_2\p_1K.\label{eq:R12def}
\end{align}
The factor \(-2\) in \(L_{12}\) occurs on both sides and cancels in the
third identity.  All multiindex binomial coefficients above are one.

From \eqref{eq:LK},
\[
\p_1K=\begin{pmatrix}0\\y_3\\-y_2/2\end{pmatrix},\qquad
\p_2K=\begin{pmatrix}y_3\\0\\-y_1/2\end{pmatrix},\qquad
\p_3K=\begin{pmatrix}y_2\\y_1\\0\end{pmatrix}.
\]
Substitution into \eqref{eq:R23def} and \eqref{eq:R13def} gives
\begin{align}
R_{23}(y)&=\left(H_{22}-\tfrac12H_{33}\right)y_1+H_{21}y_2+H_{31}y_3,
\label{eq:R23}\\
R_{13}(y)&=H_{12}y_1+\left(H_{11}-\tfrac12H_{33}\right)y_2+H_{32}y_3.
\label{eq:R13}
\end{align}

Mixed distributional derivatives commute.  Hence the first two identities
must satisfy
\[
\p_1R_{23}=\p_1\p_{23}Q_x
=\p_2\p_{13}Q_x=\p_2R_{13}.
\]
Using \eqref{eq:R23}--\eqref{eq:R13}, this necessary condition is
\[
H_{22}-\tfrac12H_{33}=H_{11}-\tfrac12H_{33},
\quad\text{or simply}\quad H_{22}=H_{11}.
\]

On the other hand, direct differentiation of \(\Phi\) gives
\begin{equation}\label{eq:Hessian}
H_{ij}(x)=\frac{3x_ix_j-|x|^2\delta_{ij}}{4\pi|x|^5}.
\end{equation}
At \(x=(1,0,0)\), formula \eqref{eq:Hessian} gives
\[
H_{11}=\frac{2}{4\pi},\qquad H_{22}=-\frac{1}{4\pi},
\]
contradicting the necessary condition.  Therefore no such factorization
exists.
\end{proof}

\section{Verification notes}

\paragraph{The obstruction is independent of polynomiality.}
Only the second derivatives appearing in Weak C6 are required.  Once they
equal the smooth linear functions \(R_{23}\) and \(R_{13}\), differentiating
in the distributional sense is legitimate, so no unmentioned \(C^3\)
hypothesis on \(P\) is needed.

\paragraph{The proposed cubic also fails directly.}
The exact symbolic checker expands the source's proposed
\(Q_x=g_1K_1+g_2K_2+g_3K_3\).  Its residuals in
\eqref{eq:R23def}--\eqref{eq:R12def} are respectively
\[
H_{11}y_1,\qquad H_{22}y_2,\qquad -\tfrac12H_{33}y_3.
\]
This calculation diagnoses the proposed route, while Theorem~\ref{thm:main}
rules out all other routes for the same \(G,L,K\).

\paragraph{The suggested differential identity is false.}
The source asks whether one can use
\(\p_{x_3}G=-\p_{x_1}G-\p_{x_2}G\).  At \(x=(1,0,0)\), the three rows of
\(D^2\Phi\) show immediately that this identity does not hold.  The proof of
Theorem~\ref{thm:main} does not depend on this observation.

\paragraph{Computational audit.}
Run
\begin{verbatim}
conda run --no-capture-output -n sandbox python \
  code/check_weak_c6.py
\end{verbatim}
from this packet directory.  The script uses exact symbolic arithmetic to
expand all three prescribed right-hand sides, check the mixed-derivative
obstruction, show that a general homogeneous cubic is inconsistent at the
chosen point, and audit the paper's tentative cubic.  The computation is not
a substitute for the proof.

\section{Scope, novelty check, and review recommendation}

The theorem fully answers the exact displayed Question 6.4 in the negative.
It does not settle the broader question in Section 6 of whether ellipticity
and cancellation imply the existence of \emph{some} Weak C6 representation,
nor does it exclude other choices of \(L\) and \(K\) for \(\nabla\).

Before promotion, the run's lightweight registry, solution, and attempt
indexes were searched using arXiv:2412.09987, the exact phrase ``Weak C6'',
``WTP'', and core gradient-representation terms.  A bounded arXiv/web search
used the exact question phrase, the source title and id, and close Weak C6
terms.  It found the source paper, whose current v2 is dated 15 May 2025 and
still poses Question 6.4, and no separate answer.  This records the search
bounds and is not an exhaustive novelty claim.

Human review is recommended.  The two most important checks are (i) the
componentwise extraction of \eqref{eq:R23def}--\eqref{eq:R12def} from Weak C6,
including the composition order and the \(-2\) coefficient, and (ii) the
mixed-derivative compatibility step.  Both are explicit and local.

\begin{thebibliography}{9}
\bibitem{Zhang}
W. Q. Zhang,
\emph{Stein-Weiss, and power weight Korn type Hardy-Sobolev Inequalities in
\(L^1\) norm}, arXiv:2412.09987v2, 2025.
\url{https://arxiv.org/abs/2412.09987}.
\end{thebibliography}

\end{document}
